\documentclass[11pt]{article}
\usepackage[margin=1in]{geometry}
\usepackage{amsmath,amssymb}

\title{Math Behind the Continuous JetBot Go-To Controller}
\author{}
\date{\today}

\begin{document}
\maketitle

\section{Purpose}

This note explains the control law implemented in
\texttt{scripts/vicon\_goto\_continuous\_viewer.py}. The controller uses the
latest Vicon pose of the JetBot and continuously computes left and right motor
commands that steer the robot toward a target point in the plane.

\section{State and Target}

Let the robot position in the Vicon plane be
\[
\mathbf{p} =
\begin{bmatrix}
x \\
y
\end{bmatrix},
\]
and let the target position be
\[
\mathbf{p}_d =
\begin{bmatrix}
x_d \\
y_d
\end{bmatrix}.
\]

The displacement to the target is
\[
\Delta \mathbf{p} =
\mathbf{p}_d - \mathbf{p}
=
\begin{bmatrix}
\Delta x \\
\Delta y
\end{bmatrix}
=
\begin{bmatrix}
x_d - x \\
y_d - y
\end{bmatrix}.
\]

The distance to the goal is
\[
r = \lVert \Delta \mathbf{p} \rVert_2
= \sqrt{(\Delta x)^2 + (\Delta y)^2}.
\]

This corresponds to the code lines that compute
\texttt{dx}, \texttt{dy}, and \texttt{distance}.

\section{Heading Estimate}

The Vicon rigid body provides a full 3D orientation. The controller extracts
the robot's physical forward axis from the rotated body-$Y$ direction and then
uses only its projection onto the ground plane.

Let $\theta$ denote the robot heading in the horizontal plane. The script
computes
\[
\theta = \operatorname{atan2}(f_y, f_x) + \delta,
\]
where:
\begin{itemize}
\item $(f_x, f_y)$ is the horizontal forward direction from the Vicon pose,
\item $\delta$ is the user-entered angle correction in radians.
\end{itemize}

The desired target heading is
\[
\theta_d = \operatorname{atan2}(\Delta y, \Delta x).
\]

The heading error is then wrapped to $(-\pi, \pi]$:
\[
e_\theta = \operatorname{wrapToPi}(\theta_d - \theta).
\]

\section{Arrival Condition}

If the robot is already close enough to the target, the controller stops:
\[
\text{if } r \le \varepsilon, \quad u_L = 0,\; u_R = 0,
\]
where $\varepsilon$ is the user-selected goal tolerance.

\section{Forward-Speed Schedule}

The controller uses a base forward speed that depends on the target distance.
Let:
\begin{itemize}
\item $v_{\max}$ be the nominal forward speed (\texttt{self.\_speed}),
\item $v_{\min}$ be the minimum forward speed near the goal
(\texttt{self.\_min\_forward\_speed}),
\item $r_s$ be the slow radius (\texttt{self.\_slow\_radius}).
\end{itemize}

Then the base forward term is
\[
v_{\text{base}} =
\begin{cases}
v_{\max}, & r \ge r_s, \\
v_{\min} + (v_{\max} - v_{\min}) \dfrac{r}{r_s}, & r < r_s.
\end{cases}
\]

So outside the slow radius the controller uses full commanded forward speed,
and inside the slow radius it linearly ramps down toward $v_{\min}$.

\section{Heading-Alignment Scaling}

The controller reduces forward motion when the robot is pointed away from the
target. It computes
\[
a = \max(0, \cos e_\theta).
\]

Then the actual forward command becomes
\[
v = v_{\text{base}} a^2.
\]

This has a useful effect:
\begin{itemize}
\item if $e_\theta \approx 0$, then $\cos e_\theta \approx 1$ and $v$ stays large,
\item if the robot is moderately misaligned, $v$ is reduced,
\item if the robot points more than $90^\circ$ away from the target, then
$\cos e_\theta < 0$, so $a = 0$ and the forward term becomes zero.
\end{itemize}

\section{Steering Term}

The steering command is proportional to heading error:
\[
\omega = \operatorname{sat}\!\left(k_\theta e_\theta,\; -u_{\text{turn}},\; u_{\text{turn}}\right),
\]
where:
\begin{itemize}
\item $k_\theta$ is the heading gain (\texttt{self.\_heading\_gain}),
\item $u_{\text{turn}}$ is the turn-speed limit (\texttt{self.\_turn\_speed}),
\item $\operatorname{sat}(x, a, b)$ clamps $x$ into the interval $[a,b]$.
\end{itemize}

\section{Differential-Drive Motor Commands}

Finally, the left and right motor commands are
\[
u_L = \operatorname{sat}(v - \omega,\; -1,\; 1),
\qquad
u_R = \operatorname{sat}(v + \omega,\; -1,\; 1).
\]

This is the core differential-drive mixing rule:
\begin{itemize}
\item the shared term $v$ pushes both wheels forward,
\item the signed steering term $\omega$ speeds one wheel up and slows the other down,
\item clamping keeps the commands inside the valid motor range.
\end{itemize}

\section{Controller Interpretation}

Putting everything together, each control cycle does the following:
\begin{enumerate}
\item read the latest Vicon pose,
\item compute the target displacement and distance,
\item compute current heading and heading error,
\item stop if the robot is inside the goal tolerance,
\item choose a base forward speed from the distance,
\item attenuate forward motion if the heading is poor,
\item compute a bounded steering correction,
\item convert forward plus steering into left and right wheel speeds.
\end{enumerate}

In compact form,
\begin{align*}
r &= \sqrt{(\Delta x)^2 + (\Delta y)^2}, \\
\theta_d &= \operatorname{atan2}(\Delta y, \Delta x), \\
e_\theta &= \operatorname{wrapToPi}(\theta_d - \theta), \\
v_{\text{base}} &=
\begin{cases}
v_{\max}, & r \ge r_s, \\
v_{\min} + (v_{\max} - v_{\min}) \dfrac{r}{r_s}, & r < r_s,
\end{cases} \\
v &= v_{\text{base}} \max(0, \cos e_\theta)^2, \\
\omega &= \operatorname{sat}(k_\theta e_\theta,\; -u_{\text{turn}},\; u_{\text{turn}}), \\
u_L &= \operatorname{sat}(v - \omega,\; -1,\; 1), \\
u_R &= \operatorname{sat}(v + \omega,\; -1,\; 1).
\end{align*}

\section{Connection to the Code}

The most relevant code block is in
\texttt{ContinuousGoToController.run()}, especially the section around:
\begin{itemize}
\item pose and distance computation,
\item heading and heading-error computation,
\item slow-radius logic,
\item forward/steering mixing,
\item left/right command output.
\end{itemize}

\end{document}
